\begin{lemma}
In a $3$-uniform $3$-simple local extremal setup $S$ at scale $D$,
\[
\sum_{g\in V_2}\sigma(g)\ge2w(X_b)-D|V_1|.
\]
\end{lemma}
\begin{proof}
Use the set, weight, cross-degree and column formulas in
\noderef{ThreeUniformThreeSimpleLocalSetup}.
By \noderef{SaturatedNeighborClassPartition}, every neighbor index of
$f$ belongs to a unique class $C_u=\{g\ne f:u\in g\}$, $u\in f$.
For $x\in X_b$ write $a_{u,x}=\deg_{\mathcal F}(u,x)$.
Since $x\notin f$, this is precisely the number of indices in $C_u$
containing $x$. The weight formula shows that $2w(X_b)$ counts tuples
$(x,u,v,g,h)$ with $x\in X_b$, distinct $u,v\in f$,
$g\in C_u$, $h\in C_v$, and $x\in g\cap h$.
Indeed the two orientations of each two-element row set give the
ordered row pairs, and independent choices of the two indices give
the product $a_{u,x}a_{v,x}$.

Partition these tuples by the first index $g$, which lies in
$N=V_1\mathbin{\dot\cup}V_2$. If $g\in V_2$, its two outside vertices
are exactly $g\cap X_b$, and its unique row is $u$; consequently the
number of these tuples with first index $g$ is exactly $\sigma(g)$.
If $g\in V_1$, it contains its unique vertex in $f$ and at least one
vertex in $X_s$. The sets $f,X_s,X_b$ are pairwise disjoint: both
outside sets lie in $X$, which avoids $f$, and the setup makes the
outside sets disjoint. Since $g$ has only three vertices, it has at
most one vertex in $X_b$. If it has none, there are no tuples.
Otherwise let that vertex be $x$. For its fixed row $u$, the possible
second indices number $\sum_{v\in f\setminus\{u\}}a_{v,x}$, at most
$\sum_{v\in f}a_{v,x}\le D$ by the column bound and nonnegativity of
the entries. Thus the total contribution from first indices in $V_1$
is at most $D|V_1|$. Subtracting this contribution from $2w(X_b)$
proves the claim. Throughout, choices are of indices, not just distinct
vertex sets.
\end{proof}
